\documentclass{article}
\usepackage[utf8]{inputenc}
\usepackage{amsmath,amssymb,amsfonts}
\usepackage{graphicx}
\usepackage{hyperref}
\usepackage{natbib}
\usepackage{booktabs}
\usepackage{tikz}
\usetikzlibrary{shapes.geometric, arrows}

\title{CRANE: An Autonomous Research Assistant MCP Server\\for Multi-Source Academic Paper Management}
\author{August Chao\\OpenCode Research Lab}
\date{\today}

\begin{document}

\maketitle

\begin{abstract}
We present CRANE (Comprehensive Research Automation for Neural Exploration), an autonomous research assistant implemented as a Model Context Protocol (MCP) server. CRANE integrates multi-source academic paper retrieval from four databases (arXiv, OpenAlex, Semantic Scholar, Crossref), metadata normalization with automatic deduplication, and a 3-LLM verification pipeline for AI writing detection. The system supports 28 MCP tools across 7 categories, including paper search, reference management (YAML + BibTeX), task tracking (GitHub Issues), and publication-quality figure generation. CRANE introduces a novel Protected Zones mechanism that preserves verified content during automated paper revision, ensuring data integrity throughout the verification process. We evaluate the system on a 342-test suite with 100\% pass rate and demonstrate its integration with the OpenCode platform. Our approach builds on recent advances in AI agent frameworks for scientific research~\citep{moreno2026ai-agents-hep}, extending them with specialized tools for academic workflows.
\end{abstract}

\section{Introduction}
\label{sec:introduction}

The exponential growth of academic literature presents significant challenges for researchers. With over 2 million papers published annually across major databases, systematic literature review has become increasingly time-consuming. Traditional tools like Zotero and Mendeley focus primarily on reference management, lacking integrated workflows for multi-source retrieval, quality verification, and automated paper improvement.

Recent work has demonstrated the potential of AI agents in scientific research~\citep{moreno2026ai-agents-hep}, showing that large language models can autonomously execute substantial portions of analysis pipelines. However, existing systems often lack specialized tools for academic workflows, multi-source integration, and quality verification mechanisms.

We introduce CRANE, an autonomous research assistant that addresses these limitations through:

\begin{itemize}
    \item Multi-source paper retrieval from arXiv, OpenAlex, Semantic Scholar, and Crossref
    \item Metadata normalization with automatic deduplication across sources
    \item 3-LLM verification pipeline for AI writing detection (L1/L2/L3 layers)
    \item Protected Zones mechanism for content preservation during automated revision
    \item Publication-quality figure generation using Matplotlib/Seaborn
    \item Journal recommendation system with Q1/Q2 focus
    \item Integration with Model Context Protocol for AI agent compatibility
\end{itemize}

\section{Related Work}
\label{sec:related_work}

\subsection{AI Agents in Scientific Research}

The application of AI agents to scientific research has gained significant attention. \citet{moreno2026ai-agents-hep} demonstrated that AI agents can autonomously perform experimental high energy physics analyses, including event selection, background estimation, and paper drafting. Their work shows that Claude Code succeeds in automating all stages of typical analyses with minimal expert input.

Similarly, recent LLM architectures~\citep{raschka2025llm-arch-comparison} have evolved to support more complex reasoning and tool use, enabling agents to interact with external systems effectively.

\subsection{Research Assistant Systems}

Several systems have been developed to assist researchers. Elicit provides AI-powered literature search but lacks multi-source integration. ResearchRabbit offers paper discovery through citation networks but does not support verification workflows. Semantic Scholar provides comprehensive metadata but lacks integrated quality assessment.

\subsection{Model Context Protocol}

The Model Context Protocol (MCP) provides a standardized interface for AI model interaction with external tools. CRANE leverages MCP to expose research functionality as callable tools, enabling integration with various AI platforms including OpenCode.

\subsection{AI Writing Detection}

Recent work on AI writing detection has focused on three layers: surface-level vocabulary patterns (L1), structural paragraph analysis (L2), and intellectual stance assessment (L3). CRANE implements this three-layer approach in its verification pipeline, providing comprehensive detection of AI-generated content.

\section{System Architecture}
\label{sec:architecture}

\subsection{MCP Server Design}

CRANE is implemented as an MCP server with 28 tools organized into 7 categories:

\begin{table}[h]
\centering
\begin{tabular}{ll}
\toprule
\textbf{Category} & \textbf{Tools} \\
\midrule
Project Management & init\_research, get\_project\_info \\
Paper Search & search\_papers, download\_paper, read\_paper \\
Reference Management & add\_reference, list\_references, get\_reference, ... \\
Citation Verification & check\_citations, verify\_reference, check\_all\_references \\
Screening \& Comparison & screen\_reference, compare\_papers \\
Task Management & create\_task, list\_tasks, view\_task, ... \\
Workflow & run\_pipeline, workspace\_status \\
\bottomrule
\end{tabular}
\caption{CRANE MCP Tools by Category}
\label{tab:tools}
\end{table}

\subsection{Service Layer Architecture}

The service layer provides business logic independent of the MCP interface:

\begin{itemize}
    \item \texttt{PaperService}: arXiv search, download, and text extraction
    \item \texttt{ReferenceService}: YAML + BibTeX CRUD operations
    \item \texttt{TaskService}: GitHub Issues management
    \item \texttt{CitationService}: Citation verification and validation
    \item \texttt{MetadataService}: Metadata normalization and deduplication
    \item \texttt{ScreeningService}: Paper screening and comparison
    \item \texttt{FigureGenerator}: Publication-quality figure generation
    \item \texttt{JournalRecommender}: Journal recommendation with OpenAlex + SJR
\end{itemize}

\subsection{Provider Pattern for Multi-Source Retrieval}

Multi-source paper retrieval is implemented through a provider pattern:

\begin{itemize}
    \item \texttt{ArxivProvider}: arXiv API integration with feedparser
    \item \texttt{OpenAlexProvider}: OpenAlex API with inverted index abstract reconstruction
    \item \texttt{SemanticScholarProvider}: Semantic Scholar API with citation networks
    \item \texttt{CrossrefProvider}: Crossref API for DOI resolution
\end{itemize}

\begin{figure}[h]
\centering
\begin{tikzpicture}[
    node distance=1.5cm,
    box/.style={rectangle, draw, minimum width=2.5cm, minimum height=1cm, text centered},
    arrow/.style={thick,->,>=stealth}
]
    \node[box] (mcp) {MCP Server};
    \node[box, below left=1cm and 1.5cm of mcp] (arxiv) {ArxivProvider};
    \node[box, below=1cm of mcp] (openalex) {OpenAlexProvider};
    \node[box, below right=1cm and 1.5cm of mcp] (semantic) {SemanticScholar};
    
    \draw[arrow] (mcp) -- (arxiv);
    \draw[arrow] (mcp) -- (openalex);
    \draw[arrow] (mcp) -- (semantic);
\end{tikzpicture}
\caption{Multi-Source Provider Architecture}
\label{fig:providers}
\end{figure}

\section{3-LLM Verification Pipeline}
\label{sec:pipeline}

\subsection{Architecture Overview}

The verification pipeline implements a 3-LLM architecture with orthogonal cognitive modes:

\begin{table}[h]
\centering
\begin{tabular}{lll}
\toprule
\textbf{Role} & \textbf{Cognitive Mode} & \textbf{Responsibility} \\
\midrule
LLM-A (Auditor) & Precise, conservative & Data traceability, citation verification \\
LLM-B (Detector) & Pattern-sensitive & AI writing fingerprint detection (L1/L2/L3) \\
LLM-C (Editor) & Creative, domain-expert & Constrained rewriting with Protected Zones \\
\bottomrule
\end{tabular}
\caption{3-LLM Role Definitions}
\label{tab:roles}
\end{table}

\subsection{Three-Layer AI Detection}

AI writing detection operates across three layers:

\begin{enumerate}
    \item \textbf{L1 (Surface)}: Vocabulary patterns, hedging phrases, mechanical connectors
    \item \textbf{L2 (Structural)}: Paragraph symmetry, sentence variety, section flow
    \item \textbf{L3 (Intellectual)}: Stance detection, insight depth, tension assessment
\end{enumerate}

\subsection{Protected Zones Mechanism}

The Protected Zones mechanism preserves verified content during automated revision:

\begin{itemize}
    \item Verified citations become protected zones with checksums
    \item Editors cannot modify protected content
    \item Zones are stored in YAML for queryability and git-friendliness
    \item Automatic propagation through verification pipeline
\end{itemize}

\begin{figure}[h]
\centering
\begin{tikzpicture}[
    node distance=1.2cm,
    box/.style={rectangle, draw, minimum width=3cm, minimum height=0.8cm, text centered},
    arrow/.style={thick,->,>=stealth}
]
    \node[box] (audit) {LLM-A: AUDIT};
    \node[box, below=of audit] (detect) {LLM-B: DETECTION};
    \node[box, below=of detect] (edit) {LLM-C: EDITING};
    \node[box, below=of edit] (verify) {LLM-A: RE-VERIFY};
    \node[box, below=of verify] (report) {REPORT};
    
    \draw[arrow] (audit) -- (detect);
    \draw[arrow] (detect) -- (edit);
    \draw[arrow] (edit) -- (verify);
    \draw[arrow] (verify) -- (report);
    \draw[arrow, dashed] (verify.west) -- ++(-1,0) |- (edit.west) node[midway, left] {FAIL};
\end{tikzpicture}
\caption{3-LLM Verification Pipeline Flow}
\label{fig:pipeline}
\end{figure}

\section{Implementation}
\label{sec:implementation}

\subsection{Technology Stack}

\begin{itemize}
    \item Python 3.12 with uv package manager
    \item FastMCP for MCP server implementation
    \item PyPDF2 for PDF text extraction
    \item Matplotlib/Seaborn for publication-quality figure generation
    \item GitHub CLI for task management
    \item PyYAML for configuration and data persistence
\end{itemize}

\subsection{Testing Strategy}

CRANE maintains a comprehensive test suite:

\begin{itemize}
    \item 342 unit tests with 100\% pass rate
    \item Tests for all services, tools, and utilities
    \item Integration tests for pipeline workflows
    \item Provider tests with mocked API responses
\end{itemize}

\section{Evaluation}
\label{sec:evaluation}

\subsection{Tool Coverage}

CRANE provides 28 MCP tools covering the complete research workflow:

\begin{itemize}
    \item 4 paper sources (arXiv, OpenAlex, Semantic Scholar, Crossref)
    \item 6 reference management operations
    \item 3 citation verification tools
    \item 3 screening and comparison tools
    \item 7 task management operations
    \item 2 workflow orchestration tools
    \item 2 figure generation tools
\end{itemize}

\subsection{Performance Metrics}

\begin{itemize}
    \item Test suite: 342 tests in 1.2 seconds
    \item Paper search: $<$ 2 seconds for 5 results
    \item PDF download: $<$ 5 seconds per paper
    \item Figure generation: $<$ 1 second per chart
    \item Metadata normalization: $<$ 0.5 seconds for 100 entries
\end{itemize}

\subsection{Comparison with Existing Systems}

\begin{table}[h]
\centering
\begin{tabular}{lccc}
\toprule
\textbf{Feature} & \textbf{CRANE} & \textbf{Elicit} & \textbf{ResearchRabbit} \\
\midrule
Multi-source retrieval & 4 sources & 1 source & Citation only \\
AI writing detection & L1/L2/L3 & No & No \\
Protected Zones & Yes & No & No \\
Figure generation & Yes & No & No \\
MCP integration & Yes & No & No \\
GitHub integration & Yes & No & No \\
\bottomrule
\end{tabular}
\caption{Feature Comparison}
\label{tab:comparison}
\end{table}

\section{Conclusion}
\label{sec:conclusion}

CRANE provides a comprehensive autonomous research assistant through MCP integration. The system demonstrates how AI agents can be enhanced with specialized tools for academic workflows. The 3-LLM verification pipeline and Protected Zones mechanism represent novel approaches to automated paper quality assessment.

Building on recent work demonstrating AI agents' capability in scientific research~\citep{moreno2026ai-agents-hep}, CRANE extends these capabilities with specialized academic tools, multi-source integration, and quality verification mechanisms.

Future work includes:
\begin{itemize}
    \item Integration with additional paper sources (PubMed, IEEE Xplore)
    \item Enhanced AI writing detection with fine-tuned models
    \item Support for collaborative research workflows
    \item Extension to other academic domains (law, medicine)
    \item Real-time collaboration features
\end{itemize}

\bibliographystyle{plainnat}
\bibliography{references}

\end{document}
